\scanpage{19}
\mspage{9}
One sees \inserted{then} that one has the notion of a map $X\longrightarrow Y$ between tame spaces, which is \uline{tame in a neighborhood of a point}, and one finds

\uline{Proposition} For a map from a tame space $X$ to an ind-tame space to be tame, it is necessary s. that it be so in a neighborhood of each point.

Thus the tame maps from $X$ to $Y$ form the sections of a sheaf --- the sheaf of tame maps from $X$ to $Y$. The \inserted{tame} structure of $X$ is known when one knows the set of tame maps from $X$ to $R$ (this \unsure{holds|is necessary} in the purely set-theoretic context, with $(M1)_0$\unsure{;|?} $(M2)$ $(M3)$ $(M4)$ only), thus also when one knows the sheaf of germs of tame functions.
\mspage{18}
